\section{Runtime Monitoring、Least Privilege 与分层隔离}

课程强调了运行时监控的重要性：系统应持续监测 Agent 行为轨迹、工具调用参数、跨域访问模式和异常状态转移。一旦触发风险阈值，应能自动降权、暂停、回滚或切换人工审批。

\begin{knowledgebox}{监控可观测信号}
\begin{itemize}
    \item 行为层：异常工具序列、异常频率、越权调用尝试；
    \item 数据层：敏感字段外流、上下文污染传播；
    \item 执行层：命令执行失败重试异常、可疑外联、权限拒绝激增；
    \item 策略层：策略冲突、动态策略频繁震荡。
\end{itemize}
\end{knowledgebox}

讲者也讨论了 least privilege 与组件分层隔离：让低权限组件即使被攻破也无法直接执行高权限动作。该策略本质上是 ``限制攻击后果半径''，使系统从 ``完全失守'' 退化为 ``局部可控失效''。

\begin{warningbox}{权限设计中的高频错误}
\begin{enumerate}
    \item 把 ``开发便利'' 当作 ``生产默认''，给 Agent 过宽权限。
    \item 缺少权限衰减机制，临时权限长期残留。
    \item 缺少跨工具统一身份上下文，导致审计断链。
\end{enumerate}
\end{warningbox}

\subsection{本章小结}
监控和最小权限不是可选强化项，而是 Agent 系统的基础结构。没有运行时约束，前置防御很容易在复杂场景中被绕过。

